\documentclass[12pt,fleqn]{article}

\usepackage[margin=1in]{geometry}
\usepackage{amsmath,amssymb}
\setlength{\mathindent}{0pt}
\mathcode`x="0078
\mathcode`y="0079

\begin{document}
For every location j, from location 1 through location 8, let the corresponding
x decision variable equal 1 if a policer officer is assigned there and 0
otherwise. Since the goal is to mimimize recruiting costs by hiring as few
officers as possible, the objective function would then be
\[
\min \quad Z=\sum_{j=1}^{8}x_j.
\]
An officer protects a location when the distance to that location is
strictly less than 100 metres. The coverage constraints are:
\begin{verbatim}
x1 + x2 + x3 + x4 >= 1                  (location 1)
x1 + x2 + x3 + x4 + x5 >= 1             (location 2)
x1 + x2 + x3 + x5 + x6 + x8 >= 1        (location 3)
x1 + x2 + x4 + x7 >= 1                  (location 4)
x2 + x3 + x5 + x6 + x8 >= 1             (location 5)
x3 + x5 + x6 + x7 + x8 >= 1             (location 6)
x4 + x6 + x7 + x8 >= 1                  (location 7)
x3 + x5 + x6 + x7 + x8 >= 1             (location 8)
\end{verbatim}

So each decision var is restricted in such a way so that its value is either equal
to 0 or 1. 1 means that the officer will be assigned to that location, whereas 0
means that no officer will be assigned to that location. 

The eight coverage constraints require every important location to have at
least one officer stationed less than 100 metres away. Among all assignments
that provid this coverage, the objective function chooses the one that hires
the smallest possible ammount of officers and therefore minimizes the total
recruiting cost.

\end{document}
